% =============================================================================
% Project : THE_PHILOSOPHY_OF_VIABILITY
% Book    : Book I — Ontology of Viability
% Book ID : VS.B01
%
% File    : 4_ASSETS/diagrams/VS.B01.P01.C03.S07.SS11.P003.DG01.scope_of_unity_and_singularity_claims.tex
%
% Asset:
%   Diagram DG01 — Scope of Unity and Singularity Claims
%
% Purpose :
%   To distinguish the independent questions involved in a claim
%   that something is one, and to prevent numerical singularity,
%   ontological unity, uniqueness, and identity from being treated
%   as equivalent conclusions.
%
% =============================================================================

\begin{figure}[htbp]
\centering

\begin{tikzpicture}[
    claim/.style={
        draw,
        rounded corners,
        align=center,
        text width=0.55\textwidth,
        minimum height=0.8cm,
        inner sep=5pt,
        font=\small
    },
    specification/.style={
        draw,
        rounded corners,
        align=center,
        text width=0.68\textwidth,
        minimum height=1.0cm,
        inner sep=6pt,
        font=\small
    },
    distinction/.style={
        draw,
        rounded corners,
        align=center,
        text width=2.7cm,
        minimum height=1.25cm,
        inner sep=5pt,
        font=\small
    },
    connection/.style={
        thick,
        dashed,
        -{Stealth[length=2mm]}
    }
]

\node[claim] (initial) at (0,3.3) {
    \textbf{Claim: ``X is one''}
};

\node[specification] (specify) at (0,1.8) {
    \textbf{Specify the claim}\\
    Subject; domain; respect; criterion; scope
};

\node[distinction] (number) at (-4.8,-0.3) {
    \textbf{Numerical Singularity}\\
    One unit under a specified counting criterion
};

\node[distinction] (unity) at (-1.6,-0.3) {
    \textbf{Ontological Unity}\\
    Unity attributed to a subject in a specified respect
};

\node[distinction] (unique) at (1.6,-0.3) {
    \textbf{Uniqueness}\\
    No other item meets the relevant criterion
};

\node[distinction] (identity) at (4.8,-0.3) {
    \textbf{Identity Question}\\
    Unity or plurality of identities requires independent grounds
};

\draw[connection] (initial) -- (specify);
\draw[connection] (specify.south west) -- (number.north);
\draw[connection] (specify.south west) -- (unity.north);
\draw[connection] (specify.south east) -- (unique.north);
\draw[connection] (specify.south east) -- (identity.north);

\node[
    align=center,
    text width=0.9\textwidth,
    font=\footnotesize
] at (0,-2.0) {
    These are distinct analytical questions. The connections
    indicate clarification, not entailment. No branch establishes
    another without an additional ontological argument.
};

\end{tikzpicture}

\caption{The Analytical Scope of Unity and Singularity Claims}
\label{fig:scope-of-unity-and-singularity-claims}
\end{figure}

% =============================================================================
% END OF DIAGRAM DG01
% =============================================================================